\documentclass[conference]{IEEEtran}

\usepackage{graphicx}
\usepackage{amsmath}
\usepackage{amsfonts}
\usepackage{booktabs}
\usepackage{hyperref}
\usepackage{xcolor}
\usepackage{algorithm}
\usepackage{algorithmic}
\usepackage{cite}

\makeatletter
\def\section{\@startsection{section}{1}{\z@}
  {-2.0ex plus -0.5ex minus -0.2ex}
  {1.0ex plus 0.2ex}
  {\normalfont\large\bfseries\raggedright}}
\def\subsection{\@startsection{subsection}{2}{\z@}
  {-1.5ex plus -0.5ex minus -0.2ex}
  {0.8ex plus 0.2ex}
  {\normalfont\normalsize\bfseries\raggedright}}
\makeatother

\title{Event-Guided Motion Compensation for LiDAR\\
Temporal Distortion in Autonomous Driving}

\author{
\IEEEauthorblockN{Sahal Ansar}
\IEEEauthorblockA{Department of Engineering\\
Email: your-email@example.com}
}

\begin{document}

\maketitle

% ============================================================
\begin{abstract}
Rotating LiDAR sensors acquire 3D scene measurements sequentially
over a full rotational sweep, producing temporally non-uniform point
clouds in which each point corresponds to a distinct acquisition
timestamp. When such scans are projected into a single camera frame
using static extrinsic calibration, the implicit assumption of
temporal uniformity is violated, causing systematic spatial
misalignment between the point cloud and image features---particularly
under ego-motion and in the presence of dynamic objects. We propose a
lightweight, deterministic image-space framework that corrects this
distortion without modifying the 3D point cloud or altering sensor
calibration. Dense optical flow between consecutive RGB frames
provides per-pixel motion estimates, while a simulated event-based
confidence mask---derived from log-intensity frame differences---gates
the correction to regions of genuine temporal activity, suppressing
noise propagation in stationary areas. A recursive temporal smoothing
stage further stabilizes flow estimates across frames. The method
operates entirely in 2D image space, requires no deep learning, and
runs efficiently on standard CPU hardware. Evaluated on four KITTI
raw sequences spanning diverse driving scenarios, the proposed
approach achieves a mean reduction in LiDAR--image temporal
inconsistency of 70.3\% compared to uncorrected projection, and
outperforms naive RGB-only flow correction by a substantial margin.
Results are consistent across all tested sequences, demonstrating
the robustness of the event-gating mechanism.
\end{abstract}

% ============================================================
\section{Introduction}

Accurate spatial alignment between LiDAR point clouds and camera
images is a foundational requirement in autonomous driving
perception. Fusing metric depth from LiDAR with dense photometric
information from cameras enables robust object detection, semantic
segmentation, and environment mapping. However, a systematic and
frequently overlooked source of misalignment arises not from
calibration error but from the fundamental sensing mechanism of
rotating LiDAR sensors.

Unlike global-shutter cameras, which capture an entire image
instantaneously, a spinning LiDAR acquires measurements
sequentially. A standard 64-beam Velodyne HDL-64E, for instance,
completes one full rotation at 10~Hz, meaning each scan spans
approximately 100~milliseconds. During this interval, the ego-vehicle
continues to move, and dynamic objects in the scene change their
positions. Consequently, a single LiDAR scan is not a spatially
consistent snapshot of the world at one instant; rather, it is a
temporally non-uniform set of observations in which early-scan
points correspond to earlier vehicle poses and late-scan points
correspond to later ones.

Formally, a LiDAR scan $\mathcal{L}$ is represented as:
\begin{equation}
  \mathcal{L} = \{ (\mathbf{X}_i,\, t_i) \}_{i=1}^{N}
\end{equation}
where $\mathbf{X}_i \in \mathbb{R}^3$ is the 3D position of the
$i$-th point and $t_i \in [t_0,\, t_0 + T_{sweep}]$ is its
acquisition timestamp, with $T_{sweep} \approx 100$~ms for a 10~Hz
sensor.

When projecting $\mathcal{L}$ onto a camera frame captured at time
$t_c$, standard pipelines assume:
\begin{equation}
  t_i = t_c \quad \forall\, i
\end{equation}
This assumption is violated for all points where $t_i \neq t_c$,
introducing temporal distortion that manifests as edge displacement
along the direction of motion, skewing of vertical structures, and
depth-dependent spatial drift.

Existing approaches to this problem either require IMU-based
ego-motion integration~\cite{kitti}, access to precise per-point
timestamps from the LiDAR hardware, or resort to deep learning
architectures that sacrifice interpretability and computational
efficiency~\cite{delta2025,huang2021}. Software-only, classical
alternatives that operate directly on standard synchronized datasets
remain underexplored.

This paper addresses the temporal distortion problem with a
deterministic, image-space correction framework. Our key insight is
that the apparent image-plane motion of LiDAR-projected points due
to temporal distortion can be approximated by scaling the dense
optical flow between consecutive camera frames. Rather than applying
this correction uniformly---which propagates optical flow noise across
stationary regions---we introduce an event-based confidence mask that
gates correction to pixels exhibiting meaningful temporal intensity
change. This selective mechanism is inspired by the operation of
Dynamic Vision Sensors (DVS)~\cite{gallego2024}, which fire
asynchronously only where brightness changes exceed a threshold.

Our contributions are:
\begin{itemize}
  \item A statistically validated proof that LiDAR temporal
    distortion is measurable and significant on standard KITTI
    sequences, with a mean uncorrected inconsistency of 11.54~px.
  \item A lightweight image-space correction pipeline combining
    Farneback optical flow~\cite{farneback}, log-intensity event
    simulation, and recursive temporal smoothing---no deep learning,
    no hardware modification.
  \item An ablation study demonstrating that event-guided gating
    provides significantly superior correction over naive RGB-only
    flow application, reducing temporal inconsistency by a mean of
    70.3\% across four KITTI sequences.
\end{itemize}

% ============================================================
\section{Related Work}

\textbf{LiDAR Motion Distortion Correction.}
The problem of intra-scan motion distortion in LiDAR has been studied
primarily in the context of SLAM and localization, where IMU
measurements are used to de-skew point clouds by integrating
ego-motion between point timestamps~\cite{kitti}. These approaches
require either hardware-level timestamp access or tight IMU
integration, neither of which is available in standard synchronized
dataset pipelines.

\textbf{LiDAR--Camera Fusion.}
Multi-modal fusion of LiDAR and RGB cameras is extensively
studied~\cite{huang2021,delta2025}. Most contemporary methods employ
deep neural network architectures for cross-modal feature alignment,
including convolutional attention gates~\cite{huang2021} and
Transformer-based dense depth estimation~\cite{delta2025}. These
approaches achieve strong benchmark performance but sacrifice
deterministic interpretability and require GPU inference, limiting
deployment on resource-constrained embedded systems.

\textbf{Event-Based Vision.}
Dynamic Vision Sensors produce asynchronous per-pixel brightness
change events with microsecond temporal resolution~\cite{gallego2024}.
Several works have explored fusing event cameras with LiDAR for
high-frequency depth estimation~\cite{li2021,bartolomei2024}, and
for extrinsic calibration~\cite{song2018}. These methods assume
access to real event camera hardware. Our work simulates event-like
saliency from standard RGB frame differences, making the approach
applicable to existing autonomous driving datasets without additional
sensors.

\textbf{Optical Flow for Sensor Fusion.}
Dense optical flow has been used as a proxy for scene motion in
depth completion~\cite{yuan2019} and temporal upsampling tasks.
Classical Farneback polynomial expansion~\cite{farneback} provides
efficient dense flow without neural network training. We build on
this foundation, extending it with event-guided selective gating to
improve robustness.

% ============================================================
\section{System Overview}

\begin{figure}[htbp]
  \centerline{\includegraphics[width=0.95\columnwidth]{fig7_pipeline_diagram.png}}
  \caption{Proposed event-guided motion compensation pipeline. RGB
    frames drive both optical flow estimation and log-intensity event
    simulation. The event confidence mask gates motion correction
    selectively. Temporal smoothing stabilizes flow estimates across
    frames. The LiDAR projection feeds corrected point coordinates
    to the output.}
  \label{fig:pipeline}
\end{figure}

Given two consecutive RGB frames $I_t$, $I_{t+1}$ and a LiDAR scan
$\mathcal{L}_t$, the system performs five sequential stages:
\begin{enumerate}
  \item Project $\mathcal{L}_t$ into the image plane of $I_t$.
  \item Compute log-intensity difference to generate a simulated
    event confidence mask $C(x,y)$.
  \item Estimate dense optical flow $F(x,y)$ between $I_t$ and
    $I_{t+1}$, with recursive temporal smoothing.
  \item Apply event-gated motion correction to projected LiDAR
    coordinates.
  \item Output corrected projections aligned with $I_{t+1}$.
\end{enumerate}
The entire correction operates in 2D image space. No 3D point cloud
modification, no IMU data, and no learned parameters are required.

% ============================================================
\section{LiDAR Projection}

Projection of LiDAR point $\mathbf{X}_i$ into the camera image plane
follows the standard KITTI calibration formulation:

\begin{equation}
  \mathbf{x}_i = P_{rect}\, R_{rect}\, T_{velo \to cam}\,
  \tilde{\mathbf{X}}_i
  \label{eq:projection}
\end{equation}

where $\tilde{\mathbf{X}}_i$ denotes the homogeneous representation
of $\mathbf{X}_i$, $T_{velo \to cam} \in \mathbb{R}^{4 \times 4}$
is the rigid body transform from LiDAR to camera coordinates,
$R_{rect} \in \mathbb{R}^{4 \times 4}$ performs stereo
rectification, and $P_{rect} \in \mathbb{R}^{3 \times 4}$ is the
rectified camera projection matrix.

The resulting homogeneous image coordinate $\mathbf{x}_i =
(u_i, v_i, z_i)^\top$ is normalized as:
\begin{equation}
  (u_i', v_i') = \left(\frac{u_i}{z_i},\; \frac{v_i}{z_i}\right)
\end{equation}
Points with $z_i \leq 0$ (behind the camera plane) are discarded.
The composed projection matrix $\Pi = P_{rect} R_{rect} T_{velo \to
cam}$ is applied in a single matrix multiplication for numerical
stability.

% ============================================================
\section{Event-Based Temporal Saliency}

\subsection{Log-Intensity Difference}

We simulate event camera output from consecutive RGB frames by
computing the per-pixel log-intensity difference:

\begin{equation}
  \Delta I(x,y) = \log\bigl(I_{t+1}(x,y)\bigr) -
  \log\bigl(I_t(x,y)\bigr)
  \label{eq:logdiff}
\end{equation}

This formulation is motivated by the DVS hardware response
model~\cite{gallego2024}, in which each pixel independently measures
logarithmic brightness change. Using log-intensity rather than raw
intensity difference provides scale invariance to global illumination
changes and ensures sensitivity to relative rather than absolute
brightness variation.

\subsection{Confidence Mask}

Binary event activation is determined by a threshold $\theta$:

\begin{equation}
  E(x,y) =
  \begin{cases}
    1 & |\Delta I(x,y)| > \theta \\
    0 & \text{otherwise}
  \end{cases}
  \label{eq:events}
\end{equation}

To spatially localize events to genuine scene boundaries and suppress
activation over uniform moving surfaces---which are not triggered by
a real DVS---we additionally filter by spatial proximity to image
edges. Canny edge detection is applied to both $I_t$ and $I_{t+1}$,
and the union of the resulting edge maps is morphologically dilated
with a $7 \times 7$ kernel. Only candidate events falling within
this dilated edge region are retained. The final confidence mask is:

\begin{equation}
  C(x,y) = E(x,y) \cdot \mathbf{1}[\text{EdgeProximity}(x,y)]
  \label{eq:conf}
\end{equation}

A small morphological dilation ($3 \times 3$ kernel) is then applied
to $C$ to extend each event pixel to its 8-neighbors, compensating
for sub-pixel LiDAR projection misalignment. On the evaluation
sequences, this produces event densities of 7--41\% depending on
scene motion magnitude.

% ============================================================
\section{Optical Flow and Temporal Smoothing}

\subsection{Dense Flow Estimation}

Dense optical flow $F(x,y) = (\Delta u, \Delta v)$ representing
pixel displacement from $I_t$ to $I_{t+1}$ is estimated using the
Farneback polynomial expansion method~\cite{farneback}, with
parameters: pyramid scale 0.5, 3 levels, window size 21 pixels,
5 iterations, polynomial degree 7. These parameters balance
computational efficiency with accuracy at the motion scales observed
in KITTI urban sequences.

\subsection{Recursive Temporal Smoothing}

Raw per-frame optical flow estimates exhibit high-frequency temporal
noise due to photometric variation and occlusion boundaries. We apply
an exponential moving average across consecutive frames:

\begin{equation}
  F_t^{smooth} = \alpha F_t + (1-\alpha)\, F_{t-1}^{smooth}
  \label{eq:smoothing}
\end{equation}

with smoothing factor $\alpha = 0.7$. This acts as a first-order
low-pass filter, preserving genuine scene motion while attenuating
frame-to-frame estimation noise. The procedure is detailed in
Algorithm~\ref{alg:smooth}.

\begin{algorithm}[h]
\caption{Recursive Temporal Flow Smoothing}
\label{alg:smooth}
\begin{algorithmic}[1]
\REQUIRE Current flow $F_t$, previous smoothed flow
         $F_{t-1}^{smooth}$, factor $\alpha$
\IF{$t = 0$ (first frame)}
  \STATE $F_t^{smooth} \leftarrow F_t$
\ELSE
  \STATE $F_t^{smooth} \leftarrow
         \alpha F_t + (1-\alpha) F_{t-1}^{smooth}$
\ENDIF
\RETURN $F_t^{smooth}$
\end{algorithmic}
\end{algorithm}

% ============================================================
\section{Event-Guided Motion Compensation}

For each projected LiDAR point at image coordinate $(u_i, v_i)$,
the corrected position is:

\begin{equation}
  (u_i^*, v_i^*) =
  \begin{cases}
    (u_i + \alpha\,\Delta u_i,\;
     v_i + \alpha\,\Delta v_i) & C(u_i,v_i) = 1 \\[4pt]
    (u_i,\; v_i) & \text{otherwise}
  \end{cases}
  \label{eq:correction}
\end{equation}

where $\Delta u_i = \Delta u(u_i,v_i)$ and
$\Delta v_i = \Delta v(u_i,v_i)$ are the smoothed flow components
sampled at the point location, and $\alpha = 0.5$ is a temporal
offset scaling factor.

The scaling factor $\alpha = 0.5$ is physically motivated: for a
10~Hz LiDAR with a 100~ms sweep duration, the mean temporal offset
between any arbitrary point's acquisition time and the camera
exposure time is approximately 50~ms, or half the inter-frame
interval. Since optical flow estimates full inter-frame displacement,
scaling by 0.5 approximates the mean correction magnitude. This is a
first-order approximation; the true scale varies per point depending
on its angular scan position within the sweep.

An additional gating condition is applied: correction is only
performed when the local flow magnitude satisfies
$0.5 < \|\mathbf{F}\| < 30.0$~px. The lower bound filters
effectively stationary pixels where no correction is meaningful. The
upper bound rejects optical flow estimation failures at occlusion
boundaries.

\begin{figure}[htbp]
  \centerline{\includegraphics[width=0.95\columnwidth]{fig4_event_map.png}}
  \caption{Simulated event confidence mask overlaid on a grayscale
    scene from \texttt{drive\_0017\_sync}. White boundaries indicate
    ON events (positive log-intensity change); blue boundaries
    indicate OFF events (negative change). Correction is applied only
    within these regions, preserving static background areas.}
  \label{fig:events}
\end{figure}

% ============================================================
\section{Validation Methodology}

\subsection{Temporal Inconsistency Metric}

We quantify misalignment using a temporal inconsistency metric
$\epsilon$ that measures the disagreement between LiDAR point
displacement and image optical flow displacement across matched
point pairs in consecutive frames.

For each valid frame pair, projected LiDAR points are matched
between $t$ and $t+1$ using a KD-tree nearest-neighbor search.
For each matched pair $(i, j)$, the LiDAR displacement
$\mathbf{d}_{lidar}$ is the vector between projected positions, and
the image displacement $\mathbf{d}_{image}$ is the optical flow
sampled at the source point location. Temporal inconsistency is:

\begin{equation}
  \epsilon = \frac{1}{|\mathcal{M}|} \sum_{(i,j) \in \mathcal{M}}
             \bigl\| \mathbf{d}_{image}^{(i)} -
             \mathbf{d}_{lidar}^{(i,j)} \bigr\|
  \label{eq:metric}
\end{equation}

where $\mathcal{M}$ denotes the set of valid matched pairs after
outlier filtering.

\subsection{Self-Consistency Baseline}

To establish the noise floor of the metric independently of LiDAR,
we compute image self-consistency by comparing forward optical flow
with the negative of backward optical flow at each pixel. This
quantifies the inherent estimation noise of the flow algorithm and
provides a theoretical lower bound on achievable $\epsilon$.

\subsection{Experimental Setup}

Experiments are conducted on four KITTI raw synchronized sequences:
\texttt{drive\_0009\_sync}, \texttt{drive\_0005\_sync},
\texttt{drive\_0013\_sync}, and \texttt{drive\_0017\_sync},
covering urban intersections, residential streets, and open road
scenarios. All sequences use the \texttt{image\_02} camera and
Velodyne HDL-64E LiDAR data with standard KITTI calibration files.
No ground-truth per-point timestamps are used. The implementation
uses Python with OpenCV and NumPy; no GPU is required.

% ============================================================
\section{Results}

\subsection{Qualitative Results}

\begin{figure}[htbp]
  \centerline{\includegraphics[width=0.95\columnwidth]{fig1_before_correction.png}}
  \caption{LiDAR projection before correction
    (\texttt{drive\_0017\_sync}, frame~40). Depth-colored points
    overlaid on the RGB frame show the raw uncorrected projection
    with KITTI static calibration.}
  \label{fig:before}
\end{figure}

\begin{figure}[htbp]
  \centerline{\includegraphics[width=0.95\columnwidth]{fig2_after_correction.png}}
  \caption{Top: uncorrected projection on frame $t$. Bottom:
    event-guided corrected projection on frame $t+1$. Corrected
    points (full brightness) are shown alongside uncorrected points
    (dimmed) with displacement vectors indicating the applied
    correction. Object boundaries in the corrected projection align
    more closely with image features in the target frame.}
  \label{fig:after}
\end{figure}

\begin{figure}[htbp]
  \centerline{\includegraphics[width=0.95\columnwidth]{fig3_zoom_comparison.png}}
  \caption{Zoomed crop centered on a high-displacement region near
    object boundaries. Left: uncorrected projection on frame $t$.
    Right: corrected projection on frame $t+1$ with yellow arrows
    indicating displacement vectors for the highest-displaced points.
    Gray points (right panel) are uncorrected; colored points are
    event-guided corrected positions.}
  \label{fig:zoom}
\end{figure}

\subsection{Quantitative Results}

Table~\ref{tab:results} reports temporal inconsistency before and
after correction across all four evaluation sequences. The proposed
event-guided method achieves a mean reduction of 70.3\% compared to
the uncorrected baseline.

\begin{table}[htbp]
\centering
\caption{Temporal Inconsistency Reduction (pixels)}
\label{tab:results}
\begin{tabular}{lcccc}
\toprule
Sequence & Original & RGB-Only & Ours & Improv. \\
\midrule
0009\_sync & 16.06 & 15.07 & 3.64 & 77.3\% \\
0005\_sync &  9.87 &  9.71 & 3.11 & 68.5\% \\
0013\_sync & 13.14 & 12.77 & 3.51 & 73.3\% \\
0017\_sync &  7.08 &  7.08 & 2.79 & 60.6\% \\
\midrule
\textbf{Mean} & \textbf{11.54} & \textbf{11.16}
              & \textbf{3.26}  & \textbf{70.3\%} \\
\bottomrule
\end{tabular}
\vspace{6pt}
\begin{minipage}{\columnwidth}
{\footnotesize Improvement = reduction in LiDAR--image motion
disagreement (px). RGB-Only applies flow without event gating.}
\end{minipage}
\end{table}

\subsection{Ablation Study}

Fig.~\ref{fig:ablation} presents an ablation across three
configurations for all four sequences: (1) no correction baseline,
(2) RGB-only flow correction without event gating, and (3) the
proposed event-guided method. The self-consistency lower bound for
each sequence is shown as a dashed line.

\begin{figure}[htbp]
  \centerline{\includegraphics[width=0.95\columnwidth]{fig5_ablation_bar_chart.png}}
  \caption{Ablation study: temporal inconsistency per method across
    all four KITTI sequences. Dashed lines indicate the image
    self-consistency lower bound per sequence. RGB-only correction
    provides negligible improvement over the baseline, while the
    proposed event-guided method consistently reduces inconsistency
    across all sequences.}
  \label{fig:ablation}
\end{figure}

The results confirm that naive RGB-only flow correction provides
negligible improvement over the uncorrected baseline---a mean
reduction of only 3.3\% (11.54 to 11.16~px). This demonstrates
that blind application of optical flow to all projected points is
ineffective: noisy flow in textureless or stationary regions
introduces corrections that cancel or worsen alignment. In contrast,
the event-gated approach reduces inconsistency to a mean of 3.26~px,
within the range of the image self-consistency lower bound on three
of four sequences, suggesting the correction approaches the noise
floor of the optical flow estimator itself.

% ============================================================
\section{Limitations}

Several limitations of the current approach should be acknowledged.

\textbf{Simulated events.} The log-intensity difference formulation
approximates DVS behavior from standard RGB frames but does not
replicate the microsecond temporal resolution, per-pixel refractory
period, or hardware noise characteristics of real event cameras.
Event density on fast-motion frames remains higher than physical DVS
output, though the event-gating mechanism retains effectiveness
across the observed density range.

\textbf{Temporal offset approximation.} The correction scale
$\alpha = 0.5$ approximates the mean temporal offset ratio for a
10~Hz LiDAR but does not account for per-point variation in scan
angle and thus acquisition timestamp. A more precise correction
would require per-point timestamp access from the LiDAR hardware,
which is not available in standard KITTI raw synchronized data.

\textbf{Correction direction assumption.} The framework assumes the
camera exposure occurs near the end of the LiDAR sweep, so that
correcting early-scan points forward in time is appropriate. The
exact phase relationship between camera trigger and LiDAR sweep in
the KITTI capture system is not publicly documented; directional
error in this assumption would attenuate rather than eliminate the
correction benefit.

\textbf{Image-space only.} The correction is applied in 2D image
space and does not modify the underlying 3D point cloud geometry.
Downstream tasks requiring metrically consistent 3D structure would
benefit from extending the correction to 3D via depth-preserved
unprojection, which we leave as future work.

\textbf{Metric dependency.} The primary validation metric measures
agreement between LiDAR displacement and optical flow displacement.
Since the correction method uses optical flow as its signal source,
this metric is partially dependent on the correction input. We
mitigate this with the independent self-consistency baseline and
note that a secondary edge alignment metric shows directionally
consistent positive improvement on the majority of tested sequences.

% ============================================================
\section{Conclusion}

We presented a deterministic, image-space framework for correcting
temporal distortion in LiDAR projections arising from the sequential
scan acquisition of rotating sensors. The method combines dense
Farneback optical flow with a simulated event confidence mask derived
from log-intensity frame differences, applying selective motion
compensation only in temporally active image regions. Recursive
temporal smoothing further stabilizes correction across frames.

Evaluated on four KITTI raw sequences, the proposed approach reduces
the LiDAR--image temporal inconsistency metric by a mean of 70.3\%,
compared to 3.3\% for naive RGB-only correction. The method requires
no deep learning, no IMU data, no hardware modification, and no
per-point LiDAR timestamps, making it directly applicable to any
dataset providing synchronized LiDAR and camera frames with standard
calibration.

Future work will explore extending the correction to 3D space using
depth-preserved unprojection, incorporating per-scan-angle temporal
offset estimation, and validating the approach with real event camera
hardware on dedicated LiDAR--event datasets such as DSEC.

% ============================================================
\begin{thebibliography}{00}

\bibitem{kitti}
A.~Geiger, P.~Lenz, C.~Stiller, and R.~Urtasun,
``Vision meets robotics: The KITTI dataset,''
\textit{Int. J. Robot. Res.}, vol.~32, no.~11,
pp.~1231--1237, 2013.

\bibitem{farneback}
G.~Farneb{\"a}ck,
``Two-frame motion estimation based on polynomial expansion,''
in \textit{Proc. Scandinavian Conf. Image Anal.\ (SCIA)}, 2003,
pp.~363--370.

\bibitem{gallego2024}
G.~Gallego \textit{et al.},
``Secrets of event-based optical flow, depth and ego-motion
estimation by contrast maximization,''
\textit{IEEE Trans. Pattern Anal. Mach. Intell.}, vol.~46, no.~12,
pp.~10517--10539, 2024.

\bibitem{delta2025}
V.~Brebion \textit{et al.},
``DELTA: Dense depth from events and LiDAR using transformer's
attention,''
in \textit{Proc. IEEE/CVF CVPR Workshops (EventVision)}, 2025.

\bibitem{huang2021}
Z.~Huang \textit{et al.},
``Fusion-based feature attention gate component for vehicle
detection using event camera and frame camera,''
\textit{IEEE Sensors J.}, vol.~21, no.~18,
pp.~20759--20769, 2021.

\bibitem{li2021}
W.~Li \textit{et al.},
``Enhancing 3-D LiDAR point clouds with event-based camera,''
\textit{IEEE Trans. Instrum. Meas.}, vol.~70, pp.~1--9, 2021.

\bibitem{bartolomei2024}
L.~Bartolomei \textit{et al.},
``LiDAR-event stereo fusion with hallucinations,''
in \textit{Proc. IEEE/CVF CVPR Workshops (EventVision)}, 2024.

\bibitem{song2018}
H.~Song \textit{et al.},
``Event-based, 6-DOF camera tracking from photometric depth maps,''
\textit{IEEE Trans. Pattern Anal. Mach. Intell.}, vol.~40, no.~10,
pp.~2396--2407, 2018.

\bibitem{yuan2019}
W.~Yuan \textit{et al.},
``Temporal upsampling of depth maps using a hybrid camera,''
\textit{IEEE Trans. Vis. Comput. Graph.}, vol.~25, no.~3,
pp.~1591--1602, 2019.

\end{thebibliography}

\end{document}